\scanpage{29}
\mspage{28}
\uline{Scholie}\quad Ainsi, via $F\longmapsto\mathrm{Mulomb}(F)$,
$(\unsure{\mathcal F|\mathfrak F|\mathscr F},\deleted{\unsure{\ll}}\leq,\ll)$ s'interprète comme un ``sous-atelier'' de $\mathrm{Fig}(\unsure{M|\mathcal M})$.

Donc la donnée de l'atelier $(\mathcal F,\leq,\ll)$ équivaut à celle d'un ens $M$, et d'une partie
\[
 \deleted{\unsure{\Phi|\mathcal F}}\mathcal F_M\,
 \deleted{\unsure{\subset\mathcal P|\mathcal F}}\,
 \deleted{\Phi}\subset\mathrm{Fig}(M)\subset
 \unsure{\mathcal P|\mathrm P|\mathfrak P|\mathscr P}(\mathcal P(M))
 \tag*{(1.44)}
\]
\[
 \left(\text{\unsure{donc|d.---\illegible{1}}}\quad
 \mathcal F_M\,\deleted{\unsure{\Phi}\,\unsure{\mathcal F_M|\Phi_M}}
 \in\mathcal P(\mathcal P(\mathcal P(M)))\right)
\]
satisfaisant certaines conditions \inserted{(qui doivent traduire C1 ; C7)}.

Du coup, la notion de \inserted{(grande)} multiombre reprend un plus grand intérêt, et je me demande s'il ne faudrait pas une notation \deleted{\unsure{plus}} et un nom \unsure{plus|pluss} lapidaires.

Je voudrais expliciter les conditions sur cet ens de figures ensemblistes dans $M$\unsure{.|-} Je désignerais par $\mathrm{At}^{*}1$ etc les axiomes en termes de $(M,\deleted{\unsure{\Phi_M}}\mathcal F_M\in\mathcal P\mathcal P\mathcal P(\unsure{M|m}))$.
\deleted{\unsure{$(\Phi_1)$|$(\Phi_F)$|$(\Phi_{\mathcal F})$|(1.49)|(1.50)}} Il est entendu que $\deleted{\unsure{\Phi}}\mathcal F_M\subset\mathrm{Fig}(M)$, i.e. \deleted{\uline{\unsure{At|AE}$^{*}$1}} que les éléments de $\deleted{\unsure{\Phi}}\mathcal F_M$ sont des ensembles de parties de $M$ qui sont des \uline{figures} ensemblistes $F$ dans $M$. \deleted{\illegible{1}} \deleted{Ainsi} Considérons
